\documentclass[10pt,landscape,a4paper]{article}
\usepackage[T1]{fontenc}
\usepackage[utf8]{inputenc}
\usepackage[ngerman]{babel}
\usepackage{tikz}
\usetikzlibrary{shapes,positioning,arrows,fit,calc,graphs,graphs.standard,trees}
\usepackage[nosf]{kpfonts}
\usepackage{multicol}
\usepackage{wrapfig}
\usepackage[top=2mm,bottom=13mm,left=2mm,right=2mm]{geometry}
\usepackage[framemethod=tikz]{mdframed}
\usepackage{microtype}
\usepackage{minted}
\usepackage{titlesec}
\usepackage{keystroke}

\renewcommand\familydefault{\sfdefault}

\definecolor{mygreen}{HTML}{4F8D4F}
\definecolor{myblue}{HTML}{3C6A6A}
\definecolor{myyellow}{HTML}{B18663}
\definecolor{myred}{HTML}{B16363}
\definecolor{sectionbg}{HTML}{D5D5D5}

\tikzstyle{every node}=[font=\footnotesize\sffamily]
\tikzstyle{sup}=[rectangle,draw,fill=myblue!50,rounded corners=.8ex]
\tikzstyle{sub}=[rectangle,draw,fill=mygreen!70,rounded corners=.8ex]
\tikzstyle{block} = [rectangle, draw, text width=2.5cm, text centered, rounded corners, minimum height=4em,fill=blue!20]
\tikzstyle{line} = [draw,thick, -latex']
\tikzstyle{cloud} = [draw, ellipse, text width=2.5cm, text centered]
\tikzstyle{edge from parent}=[<-,thick,draw]

\let\bar\overline

\def\firstcircle{(0,0) circle (1.5cm)}
\def\secondcircle{(0:2cm) circle (1.5cm)}

\colorlet{circle edge}{myblue}
\colorlet{circle area}{myblue!5}

\tikzset{filled/.style={fill=circle area, draw=circle edge, thick},
    outline/.style={draw=circle edge, thick}}

\pgfdeclarelayer{background}
\pgfsetlayers{background,main}

\everymath\expandafter{\the\everymath \color{myblue}}
\everydisplay\expandafter{\the\everydisplay \color{myblue}}

\renewcommand{\baselinestretch}{.8}
\pagestyle{empty}

\global\mdfdefinestyle{header}{%
linecolor=gray,linewidth=1pt,%
leftmargin=0mm,rightmargin=0mm,skipbelow=0mm,skipabove=0mm,
}

\newcommand{\header}{
  \vspace{1mm}
  \begin{center}
    \includegraphics[height=09mm]{codecentric.png}
  \end{center}
}

\makeatletter

\newcommand{\colorsection}[1]{\colorbox{sectionbg}{\parbox{\textwidth-215\fboxsep}{\centering#1}}}
\titleformat{\section}[block]{\color{myblue}\sffamily\small\bfseries\filcenter}{}{1em}{\colorsection}
\titleformat{\subsection}[block]{\color{myblue}\sffamily\small\bfseries\filcenter}{}{1em}{}

\def\multi@column@out{%
   \ifnum\outputpenalty <-\@M
   \speci@ls \else
   \ifvoid\colbreak@box\else
     \mult@info\@ne{Re-adding forced
               break(s) for splitting}%
     \setbox\@cclv\vbox{%
        \unvbox\colbreak@box
        \penalty-\@Mv\unvbox\@cclv}%
   \fi
   \splittopskip\topskip
   \splitmaxdepth\maxdepth
   \dimen@\@colroom
   \divide\skip\footins\col@number
   \ifvoid\footins \else
      \leave@mult@footins
   \fi
   \let\ifshr@kingsaved\ifshr@king
   \ifvbox \@kludgeins
     \advance \dimen@ -\ht\@kludgeins
     \ifdim \wd\@kludgeins>\z@
        \shr@nkingtrue
     \fi
   \fi
   \process@cols\mult@gfirstbox{%
%%%%% START CHANGE
\ifnum\count@=\numexpr\mult@rightbox+2\relax
          \setbox\count@\vsplit\@cclv to \dimexpr \dimen@-1cm\relax
\setbox\count@\vbox to \dimen@{\vbox to 1cm{\header}\unvbox\count@\vss}%
\else
      \setbox\count@\vsplit\@cclv to \dimen@
\fi
%%%%% END CHANGE
            \set@keptmarks
            \setbox\count@
                 \vbox to\dimen@
                  {\unvbox\count@
                   \remove@discardable@items
                   \ifshr@nking\vfill\fi}%
           }%
   \setbox\mult@rightbox
       \vsplit\@cclv to\dimen@
   \set@keptmarks
   \setbox\mult@rightbox\vbox to\dimen@
          {\unvbox\mult@rightbox
           \remove@discardable@items
           \ifshr@nking\vfill\fi}%
   \let\ifshr@king\ifshr@kingsaved
   \ifvoid\@cclv \else
       \unvbox\@cclv
       \ifnum\outputpenalty=\@M
       \else
          \penalty\outputpenalty
       \fi
       \ifvoid\footins\else
         \PackageWarning{multicol}%
          {I moved some lines to
           the next page.\MessageBreak
           Footnotes on page
           \thepage\space might be wrong}%
       \fi
       \ifnum \c@tracingmulticols>\thr@@
                    \hrule\allowbreak \fi
   \fi
   \ifx\@empty\kept@firstmark
      \let\firstmark\kept@topmark
      \let\botmark\kept@topmark
   \else
      \let\firstmark\kept@firstmark
      \let\botmark\kept@botmark
   \fi
   \let\topmark\kept@topmark
   \mult@info\tw@
        {Use kept top mark:\MessageBreak
          \meaning\kept@topmark
         \MessageBreak
         Use kept first mark:\MessageBreak
          \meaning\kept@firstmark
        \MessageBreak
         Use kept bot mark:\MessageBreak
          \meaning\kept@botmark
        \MessageBreak
         Produce first mark:\MessageBreak
          \meaning\firstmark
        \MessageBreak
        Produce bot mark:\MessageBreak
          \meaning\botmark
         \@gobbletwo}%
   \setbox\@cclv\vbox{\unvbox\partial@page
                      \page@sofar}%
   \@makecol\@outputpage
     \global\let\kept@topmark\botmark
     \global\let\kept@firstmark\@empty
     \global\let\kept@botmark\@empty
     \mult@info\tw@
        {(Re)Init top mark:\MessageBreak
         \meaning\kept@topmark
         \@gobbletwo}%
   \global\@colroom\@colht
   \global \@mparbottom \z@
   \process@deferreds
   \@whilesw\if@fcolmade\fi{\@outputpage
      \global\@colroom\@colht
      \process@deferreds}%
   \mult@info\@ne
     {Colroom:\MessageBreak
      \the\@colht\space
              after float space removed
              = \the\@colroom \@gobble}%
    \set@mult@vsize \global
  \fi}

\makeatother
\setlength{\parindent}{0pt}

\begin{document}
\setminted{style=manni}

\small
\begin{multicols*}{4}
  \section{Variables}
  \begin{minted}{scala}
var x         = "mutable"
val y         = "immutable"

// Initialized once on first access
lazy val z    = "lazy"

// Pattern matching on left side
val (one,two) = ("one", '2')
\end{minted}

\section{Methods}
\begin{minted}{scala}
def id(c: Char): Char = c

def add(n: Int, m: Int): Int = n + m

// Curried version (one argument per list)
def add(n: Int)(m: Int): Int = n + m

// By-name parameters, evaluates `a` twice
def twice[A](a: => A) = { a; a }

// Repeated Parameters (Varargs)
def many(cs: Char*): Seq[Char] = cs

// Calling a varargs method with a `Seq`
scala> many(Seq('1', '2') :_*)
\end{minted}

\section{Function Types}
\begin{center}
\begin{minipage}{0.37\linewidth}
\begin{minted}{scala}
   ()      => A
 Unit      => A
    A      => B
(A, B)     => C
(A, .., V) => W
\end{minted}
\end{minipage}
\begin{minipage}{0.09\linewidth}
  \begin{minted}{text}
==
==
==
==
==
  \end{minted}
\end{minipage}
\begin{minipage}{0.50\linewidth}
\begin{minted}{scala}
Function0[A]
Function1[Unit, A]
Function1[A, B]
Function2[A, B, C]
Function22[A, .., V, W]
\end{minted}
\end{minipage}
\end{center}
\section{Types}
\begin{center}
\begin{tikzpicture}[auto,grow=right]
\tikzstyle{level 1}=[sibling distance=16em,level distance=3em]
\tikzstyle{level 2}=[sibling distance=4em,level distance=6em]
% Place nodes
\node [sup] (any) {Any}
child{node [sup] (anyval) {AnyVal}
    child{node [sub] (nums) {Numbers}}
    child{node [sub] (char) {Char}}
    child{node [sub,fill=myred!60] (char) {Your AnyVal}}
    child{node [sub] (bool) {Boolean}}
}
child{node [sup] (anyref) {AnyRef}
    child{node [sub] (colls) {Collections}}
    child{node [sub,fill=myred!60] (clss) {Your Class}}
    child{node [sub] (str) {String}}
};
\node[sup,xshift=1cm,right of=clss](null){Null};
\node[sup,xshift=4.5cm,right of=any](nothing){Nothing};
\node[sup,above of=anyref](obj){Object};
%% Draw edges
\path [line] (null) -- (str);
\path [line] (null) -- (colls);
\path [line] (null) -- (clss);
\path [line] (nothing) -- (null);
\path [line] (nothing) -- (bool);
\path [line] (nothing) -- (char);
\path [line] (nothing) -- (nums);
\path [line,dashed] (obj) -- (anyref);
\end{tikzpicture}
\end{center}
\section{Option, Either, Try}
\begin{center}
  \begin{tikzpicture}[
    level 1/.style={sibling distance=10em}]
    \node[sup]{Option[A]}
    child{node[sub] {None[Nothing]}}
    child{node[sub]{Some[A](value: A)}}
    ;
  \end{tikzpicture}
\end{center}
\begin{center}
  \begin{tikzpicture}[
    level 1/.style={sibling distance=10em}]
    \node[sup]{Either[A, B]}
    child{node[sub] {Left[A](value: A)}}
    child{node[sub]{Right[B](value: B)}}
    ;
  \end{tikzpicture}
\end{center}
\begin{center}
  \begin{tikzpicture}[
    level 1/.style={sibling distance=12em}]
    \node[sup]{Try[A]}
    child{node[sub] {Failure[A](e: Throwable)}}
    child{node[sub]{Success[A](value: A)}}
    ;
  \end{tikzpicture}
\end{center}
\subsection{Option}
\begin{itemize}
\item getOrElse
\item fold
\item ...
\end{itemize}
\subsection{Either}
\begin{itemize}
\item getOrElse
\item fold
\item ...
\end{itemize}
\subsection{Try}
\begin{itemize}
\item getOrElse
\item fold
\item ...
\end{itemize}

\section{Operators}
\begin{itemize}
\item \textit{precedence} is determined by the \textit{first} character of the
  operator's name:

  \begin{tabular}{l | l}
    lowest & letters \\
           & \mintinline{scala}/|/ \\
           & \mintinline{scala}|^| \\
           & \mintinline{scala}|&| \\
           & \mintinline{scala}|= !| \\
           & \mintinline{scala}|< >| \\
           & \mintinline{scala}|:| \\
           & \mintinline{scala}|+ -| \\
           & \mintinline{scala}|* / %| \\
    highest & other special characters

  \end{tabular}
\item \textit{associativity} is determined by the \textit{last}
  character of the operator's name.  If the last character is a
  \mintinline{scala}|:| it is right-associative, otherwise it is
  left-associative
\end{itemize}

\section{Ranges}

\begin{minted}{scala}
// Exclusive range: 1,2,3,4
1 until 5

// Inclusive range: 1,2,3,4,5
1 to 5

// Custom step size: 1,3,5
1.to(5,2)

\end{minted}

\section{Imports}

\begin{minted}{scala}
// Regular import
import scala.util.Try

// Import multiple
import scala.util.{Success, Failure}

// Import wildcard
import scala.util._

// Import under different name
import java.time.{Clock => JClock}

// Import everything except
import scala.util.{Try => _, _}

// Multiple import clauses
import scala.util.Try, scala.util.Success
\end{minted}

\section{Modifiers}
\begin{minted}{scala}
private
private[this]
private[<package>]

protected
protected[this]

override
abstract override
abstract
final
sealed
lazy
\end{minted}

\section{Loops}
\begin{minted}{scala}
// while-loop
while (<condition>) {
  <expression>
}

// do-while-loop
do {
  <expression>
} while (<condition>)

// for-loop
for (i <- 1 to 10) println(i)

// nested for-loop
for (i <- 1 to 10; j <- 1 to 10) println(i,j)

// for-comprehension
for (i <- 1 to 3; j <- 1 to i) yield i * y

// guards in for-loops and for-comprehensions
for (i <- 1 to 5 if i > 4) yield i

// use curly braces for multiline expressions
for {
  i <- 1 to 5
  j <- 1 to i
  if j % 2 == 1
} yield i * j
\end{minted}
\section{Futures}

\subsection{Timeouts}
\begin{itemize}
\item
  \begin{minted}{scala}
    import scala.concurrent.duration._
  \end{minted}
\item Await.result(future, 42.seconds)
\end{itemize}

\section{Implicits}
\begin{itemize}
\item implicit conversions
\item implicit arguments
\end{itemize}

\section{Type Parameters}

\section{Pattern Matching}
\begin{minted}{scala}
arg match {
  // Variable Patterns
  case x =>

  // Typed Patterns
  case x: String =>

  // Literal Patterns
  case 42 =>

  // Stable Identifier Patterns
  case `foo` =>

  // Constructor Patterns
  case Foo(x,y) =>

  // Tuple Patterns
  case (x,y,z) =>

  // Extractor Patterns
  // (See 'Custom Extractors')
  case NumString(x) =>

  // Pattern Sequences
  case x1 +: x2 +: xs =>

  // Pattern Alternatives
  case true | 42 | "str" =>

  // Pattern Binders
  case pair@(x,y) =>

}
\end{minted}

\section{Custom Extractors}

\begin{minted}{scala}
object Even {
  def unapply(n: Int): Boolean =
    n % 2 == 0
}
\end{minted}

\begin{minted}{text}
> 41 match { case Even() => '!' }
scala.MatchError: 41 ...
> 42 match { case Even() => "even!" }
res: String = "even!"
\end{minted}

\begin{minted}{scala}
object NumString {
  def unapply(s: String): Option[Int] =
    Try(s.toInt).toOption
}
\end{minted}

\begin{minted}{text}
> "42" match { case NumString(n) => n }
res: Int = 42
> "scala" match { case NumString(n) => n }
scala.MatchError: scala
\end{minted}

\begin{minted}{scala}
object Words {
  def unapplySeq(s: String): Option[Seq[String]] =
    Some(s.split("\\s+").to[Seq])
}
\end{minted}

\begin{minted}{text}
> "foo bar baz" match { case Words(ws) => ws }
res: Seq[String] = Vector("foo", "bar", "baz")
> "test" match { case Words("foo" +: _) => 1 }
scala.MatchError: test
\end{minted}

\section{IntelliJ}
\begin{tabular}{l l}
  Reformat & \Ctrl+\Alt+\keystroke{L} \\
  Extract Method & \Ctrl+\Alt+\keystroke{M} \\
  Rename & \Shift+\keystroke{F6} \\
  Expand selection & \Ctrl+\keystroke{W} \\
  Surround with & \Ctrl+\Alt+\keystroke{T}
\end{tabular}

\end{multicols*}
\end{document}
